% Generated by build.py from the modular sources.

% BEGIN preamble.tex
\documentclass[11pt,letterpaper]{article}
\usepackage[T1]{fontenc}
\usepackage[utf8]{inputenc}
\usepackage{lmodern}
\usepackage[margin=1in,headheight=15pt]{geometry}
\usepackage{amsmath,amssymb,amsthm,mathtools}
\usepackage{microtype,booktabs,longtable,array,enumitem,xcolor,fancyhdr}
\usepackage{hyperref}
\definecolor{inkblue}{HTML}{173B57}
\definecolor{linkblue}{HTML}{176B91}
\hypersetup{colorlinks=true,linkcolor=inkblue,citecolor=linkblue,urlcolor=linkblue,
 pdftitle={Ordered Hurwitz Resonances and Cubic Stieltjes Products},
 pdfauthor={Research continuation prepared for Vladimir Reshetnikov with OpenAI ChatGPT},
 pdfsubject={Exact identities, ordered multiple Hurwitz zeta, Stieltjes constants, Gamma and Tornheim derivatives}}
\pagestyle{fancy}\fancyhf{}
\fancyhead[L]{\small\textsc{Ordered Hurwitz resonances}}
\fancyhead[R]{\small ProveIt research continuation}
\fancyfoot[C]{\thepage}
\setlength{\parskip}{4pt}
\setlength{\parindent}{15pt}
\setlength{\emergencystretch}{2em}
\allowdisplaybreaks[2]
\numberwithin{equation}{section}
\newtheorem{theorem}{Theorem}[section]
\newtheorem{proposition}[theorem]{Proposition}
\newtheorem{lemma}[theorem]{Lemma}
\newtheorem{corollary}[theorem]{Corollary}
\theoremstyle{definition}
\newtheorem{definition}[theorem]{Definition}
\theoremstyle{remark}
\newtheorem{remark}[theorem]{Remark}
\DeclareMathOperator{\FP}{FP}
\DeclareMathOperator{\CT}{CT}
\DeclareMathOperator{\Li}{Li}
\DeclareMathOperator{\ord}{ord}
\newcommand{\eps}{\varepsilon}
\newcommand{\src}[1]{\texttt{\detokenize{#1}}}
\newcommand{\R}{\mathbb R}
\newcommand{\C}{\mathbb C}
\newcommand{\N}{\mathbb N}
\title{\color{inkblue}\bfseries Ordered Hurwitz Resonances\\[6pt]
and Cubic Stieltjes Products\\[10pt]
\large Exact Laurent coefficients, Gamma generating functions,\\
sharp curve dependence, and a polylogarithmic Mellin bridge}
\author{Research continuation prepared for Vladimir Reshetnikov\\
\small Mathematical derivations and independent review with OpenAI ChatGPT}
\date{10 October 2026}

% END preamble.tex

\begin{document}
\maketitle

% BEGIN sections/00_abstract.tex
\begin{abstract}
We answer the ordered-depth-three and tangential-approach questions in
the recent ProveIt resonant-jets programme, and obtain a stronger exact
evaluation family at every depth. For a positive common shift, an
elementary Hurwitz-tail subtraction constructs the canonical ordered
Laurent germs with absolutely convergent harmonic--Stieltjes
coefficients. Their diagonal regular generating function is explicit,
and at zero deformation it is the Gamma quotient
$\Gamma(a)/\Gamma(a+y)$. It follows that the constant Laurent coefficient
of $Z_d(1+p\eps,1+q\eps,\ldots,1+q\eps;a)$ is a finite expression
in ordinary Hurwitz derivatives and generalized Stieltjes constants,
for every depth and every transverse pair of slopes. Explicit formulas
are given through depth five. At depth three, arbitrary independent
slopes require only one additional coefficient, identified precisely
with a harmonic Stieltjes constant.

For analytic curves with prefix vanishing orders $\nu_1,\ldots,\nu_d$,
we prove an exact finite coefficient formula and show that the curve
jet through order $\sum\nu_j+\max\nu_j$ is both sufficient and, in the
worst case, necessary. A separate endpoint calculation constructs all
cubic coincident Stieltjes moments, including arbitrary argument
derivatives. The specified finite part of $\int_0^1\psi(x)^3\,dx$ is
reduced to one diagonal Tornheim third derivative, with an independent
polylogarithmic Mellin representation and numerical check.

The general ordered Laurent framework and the triple-Hurwitz/Tornheim
correspondence have established antecedents and are credited explicitly.
No arithmetic independence or first-publication claim is made. The
short Gaussian $S_6$ and current $S_8$ identities remain conjectural,
and reduction of the remaining Tornheim derivative to depth-one
constants remains open here.
\end{abstract}

% END sections/00_abstract.tex

\clearpage
\begingroup
\small
\setlength{\parskip}{0pt}
\tableofcontents
\endgroup
\clearpage

% BEGIN sections/01_scope.tex
\section{Research targets, results, and source status}
\label{sec:scope}

The question motivating this continuation is how much of a singular
\emph{ordered} nested sum can be evaluated using the familiar Gamma,
polygamma, Hurwitz-zeta and Stieltjes functions. Full symmetrization
removes ordering information and often reduces by the stuffle product.
The individual ordered germ retains that information. Its singularities
also meet along several hyperplanes, so a constant Laurent coefficient
must specify the approach to their intersection.

We work from the repository snapshot
\begin{center}
\small\texttt{4c1286161c425eca957b4b9f7cd436c37a4c98f3}.
\end{center}
The source survey includes the canonical manuscript, its status records,
and the seven ZIP packages then present in \src{docs/incoming}. The
mathematical audit concentrates on the relevant ordered-germ,
coincident-product, resonance, and Gaussian-status material; it is
not a re-verification of every proof in the full manuscript.
The exact files and archive hashes are recorded in the accompanying
\src{integration/source_snapshot.json}. All sources were inspected as
research evidence; successful numerical replay is not treated as a
proof of an analytic theorem.

\subsection{The concrete questions answered}

The incoming \emph{Resonant Gamma Jets and Directional Zeta Identities}
report asks for an ordered depth-three germ with independent exponents,
its intersecting polar divisors, a normally convergent remainder, and
the constant along every transverse ray. It separately asks for
tangential curves and the exact number of curve coefficients needed
to determine a finite part \cite[Further research, Q1 and Q3]{repojets}.
Both questions are answered here. The first answer also yields the
following all-depth specialization:
\begin{equation}\label{scope:mainfamily}
 \FP_{\eps=0}Z_d(1+p\eps,1+q\eps,\ldots,1+q\eps;a),
 \qquad a>0.
\end{equation}
It closes in ordinary Hurwitz-zeta derivatives and Stieltjes constants
whenever all prefix slopes $p+jq$ are nonzero. The initial slope may
differ from every trailing slope; no symmetrization of the ordered sum
is needed.

The \emph{Coincident Stieltjes} report asks for a cubic closure theorem
and specifically for the coordinate-normalized digamma cube
\cite[Further research: cubic coincident moments]{repocoincident}.
We construct the full cubic completion and evaluate that finite part
in terms of a single precisely defined diagonal Tornheim derivative.
This is a finite reduction to an established higher-depth function.
A further reduction of that derivative to ordinary constants is a
separate unanswered question.

\begin{center}
\begin{tabular}{>{\raggedright\arraybackslash}p{.26\textwidth}>{\raggedright\arraybackslash}p{.64\textwidth}}
\toprule
Result & Exact outcome\\
\midrule
Ordered local germs & Convergent remainder construction at every depth and
every common positive shift; finite coefficient recursion.\\[4pt]
One distinguished slope & Every nonpositive Laurent coefficient reduces
to finitely many ordinary Hurwitz and Stieltjes jets.\\[4pt]
Independent depth three & Complete pole part and constant term; one
harmonic Stieltjes coefficient contains the residual ordering dependence.\\[4pt]
Tangential curves & Pole order $\sum\nu_j$ and a sharp required curve
jet order $\sum\nu_j+\max\nu_j$.\\[4pt]
Cubic coincident products & All Stieltjes indices and all argument
derivatives; explicit digamma-cube reduction to $\tau'''(0)$.\\
\bottomrule
\end{tabular}
\end{center}

\subsection{What is classical and what is added to this programme}

Laurent expansions of Euler--Zagier multiple zeta functions have a
substantial literature. Matsumoto, Onozuka and Wakabayashi analyze Laurent
coefficients at integer points \cite{MOW}. Saha constructs regularized
multiple Stieltjes germs and, in particular, gives the canonical
ordered decomposition at the all-one point \cite[Theorem 2]{Saha}.
Our tail-defect proof realizes that structure for an arbitrary
positive common Hurwitz shift. The explicit slope evaluations, the
reduction of every transverse depth-three directional finite part to
one named harmonic coefficient and ordinary Hurwitz data, and the
sharp curve-jet statement are the
specific conclusions developed for the repository's current questions.

The harmonic Hurwitz-zeta convention is that of Karg\i n, Dil,
Cenkci and Can \cite{Kargin}; the strict ordered convention used
here differs from their inclusive sum by a single Hurwitz-zeta term.
We give the conversion rather than silently identifying the constants.
The triple-Hurwitz Fourier relation is part of the Tornheim literature
of Espinosa and Moll \cite{EMTornheim}. The Mellin method follows the
established approach discussed by Borwein and Dilcher and by Bailey
and Borwein \cite{BorweinDilcher,BaileyBorwein}. Low diagonal derivatives
are recovered and used as consistency checks, without claiming their
first evaluation. The bibliography is a targeted source comparison,
not an exhaustive proof of publication priority.

\subsection{Conventions and mathematical status}

For $a>0$ we use real logarithms of $n+a$, and
\begin{equation}\label{scope:stieltjes}
 \zeta(1+u,a)=\frac1u+
  \sum_{m\ge0}\frac{(-1)^m}{m!}\gamma_m(a)u^m,
 \qquad \gamma_0(a)=-\psi(a).
\end{equation}
A prime on $\zeta$ denotes its first, spectral derivative. For $m\ge2$,
\begin{equation}\label{scope:polygamma}
 \zeta(m,a)=\frac{(-1)^m}{(m-1)!}\psi^{(m-1)}(a).
\end{equation}
These conventions follow the standard Hurwitz calculus \cite{DLMFHurwitz}.
Without a shift, $\gamma_m=\gamma_m(1)$ and $\gamma=\gamma_0(1)$.
The rising factorial is $(z)_k=\prod_{j=0}^{k-1}(z+j)$, with
$(z)_0=1$.
Brackets $[u^j]$ or $[y^jt^k]$ mean ordinary Taylor coefficients,
so conversion to derivatives requires the displayed factorials.

For a specified meromorphic function of one variable, $\FP$ means its
coefficient of $\eps^0$. For a spatial endpoint integral, it means
the constant remaining in the cutoff expansion in the stated spatial
coordinate. These two operations are related by explicit corrections
in Section~\ref{sec:cubic}; they are not interchangeable definitions.
The holomorphic regular germ is a third object, defined by an exact
subtraction recursion.

The canonical manuscript already proves $S_4$ and retains the short
$S_6$ and revised $S_8$ candidates as conjectures \cite{repomain}.
No argument in this article changes that status. We do not infer
period independence, minimal arithmetic bases, or nonreducibility
from a failed search. Every theorem below is proved analytically;
the attached computations check finite algebra and independently
evaluated examples.

% END sections/01_scope.tex


% BEGIN sections/02_ordered.tex
\section{Ordered Hurwitz germs and a family closing in ordinary zeta jets}

Fix $a>0$. All logarithms of $n+a$ are real. Write
\[
 Z_d(s_1,\ldots,s_d;a)
 =\sum_{n_1>\cdots>n_d\ge0}\prod_{j=1}^d(n_j+a)^{-s_j},
 \qquad Z_0=1.
\]
The defining series converges normally whenever
$\Re(s_1+\cdots+s_j)>j$ for $1\le j\le d$.
Put $u_j=s_j-1$ and $P_k=u_1+\cdots+u_k$.
Empty products equal one.

\subsection{An elementary holomorphic remainder construction}

Define the tail defect
\begin{equation}\label{ord:taildefect}
 D(s,x)=\zeta(s,x+1)-\frac{x^{1-s}}{s-1},\qquad x>0.
\end{equation}
Its apparent singularity at $s=1$ is removable. On $\Re s>0$ it has
the absolutely convergent representation
\begin{equation}\label{ord:tailmellin}
 D(s,x)=\frac1{\Gamma(s)}\int_0^\infty
 t^{s-1}e^{-xt}\left(\frac1{e^t-1}-\frac1t\right)\,dt.
\end{equation}
Initially this follows by subtracting the usual Mellin integrals in
$\Re s>1$. The bracket is bounded near zero and is $O(1/t)$ at
infinity, proving continuation and normal convergence in $\Re s>0$.
In particular, uniformly for $s$ on a compact subset there,
$D(s,x)=O(x^{-\Re s})$ as $x\to\infty$.

For $d\ge2$, define
\begin{equation}\label{ord:Edef}
 E_d(u_1,\ldots,u_d;a)
 =\sum_{n_2>\cdots>n_d\ge0}
 D(1+u_1,n_2+a)\prod_{j=2}^d(n_j+a)^{-1-u_j}.
\end{equation}
This is holomorphic on the $\ell^1$ ball
\[
 \mathcal U_d=\{\mathbf u\in\mathbb C^d:
                         |u_1|+\cdots+|u_d|<1/2\}.
\]
Indeed, on a compact subset the inner finite sum over $n_3,\ldots,n_d$
is bounded by a constant times
$(1+n_2)^{\sum_{j=3}^d|u_j|}(1+\log(1+n_2))^{d-2}$.
The remaining outer summand is therefore
$O((1+n_2)^{-2+\sum|u_j|}(1+\log(1+n_2))^{d-2})$.
This is summable uniformly on compact subsets. Holomorphy follows by
normal convergence; fixed spectral derivatives can also be taken termwise.

Set $R_0=1$ and
\[
 R_1(u;a)=\zeta(1+u,a)-\frac1u.
\]
Recursively, for $d\ge2$, put
\begin{equation}\label{ord:Rrecursion}
 R_d(u_1,\ldots,u_d;a)=E_d(u_1,\ldots,u_d;a)
 +\frac{R_{d-1}(u_1+u_2,u_3,\ldots,u_d;a)
             -R_{d-1}(u_2,u_3,\ldots,u_d;a)}{u_1}.
\end{equation}
The last quotient is its holomorphic divided difference at $u_1=0$;
equivalently, it equals the integral from zero to one of the first
partial derivative of $R_{d-1}$ at
$(u_2+\tau u_1,u_3,\ldots,u_d)$. Thus every $R_d$ is holomorphic on
$\mathcal U_d$.

\begin{theorem}[Canonical shifted ordered decomposition]
\label{ord:canonical}
As meromorphic germs on $\mathcal U_d$,
\begin{equation}\label{ord:canonicalformula}
 \boxed{\displaystyle
 Z_d(1+u_1,\ldots,1+u_d;a)
 =\sum_{k=0}^{d}
 \frac{R_{d-k}(u_{k+1},\ldots,u_d;a)}{P_1P_2\cdots P_k}.}
\end{equation}
The local polar hyperplanes are precisely $P_k=0$, each generically
simple in the independent prefix coordinates.
\end{theorem}

\begin{proof}
Summing first over $n_1>n_2$ in the absolute-convergence domain gives
\[
 Z_d(1+\mathbf u;a)
 =\frac1{u_1}Z_{d-1}(1+u_1+u_2,1+u_3,\ldots,1+u_d;a)
       +E_d(\mathbf u;a).
\]
Apply the induction hypothesis to the depth $d-1$ function. Every term
except its holomorphic remainder already has the required denominator
$P_1\cdots P_k$. Splitting the remaining numerator into
$R_{d-1}(u_2,\ldots,u_d)$ plus its difference gives exactly
\eqref{ord:Rrecursion}. Normal convergence and the identity theorem then
establish the meromorphic identity throughout $\mathcal U_d$.
In the coordinates $P_1,\ldots,P_d$, every displayed denominator is
square-free. The coefficient of $P_1^{-1}\cdots P_d^{-1}$ is one,
because only the $k=d$ term contributes. Consequently no one of these
polar divisors is removable as a germ.
\end{proof}

For $a=1$, the shape of \eqref{ord:canonicalformula} is the ordered
Laurent decomposition of Saha\cite[Theorem 2]{Saha}. The construction above
provides a short proof for arbitrary common shift $a>0$ and supplies
absolutely convergent representatives for every regular coefficient.
The canonical framework is credited to that work; the explicit shifted
construction and the evaluations below are derived here.

\subsection{Explicit coefficients in harmonic--Stieltjes coordinates}

Let
\[
 Q_m(x)=\gamma_m(x+1)+\frac{\log^{m+1}x}{m+1}.
\]
In particular $Q_0(x)=\log x-\psi(x+1)$, and
\[
 D(1+u,x)=\sum_{m\ge0}\frac{(-1)^m}{m!}Q_m(x)u^m.
\]
Consequently, for a multi-index $\mathbf m=(m_1,\ldots,m_d)$,
\begin{equation}\label{ord:Ecoeff}
 [\mathbf u^{\mathbf m}]E_d(\mathbf u;a)
 =\frac{(-1)^{|\mathbf m|}}{m_1!\cdots m_d!}
 \sum_{n_2>\cdots>n_d\ge0}
 \frac{Q_{m_1}(n_2+a)\log^{m_2}(n_2+a)}{n_2+a}
 \prod_{j=3}^d\frac{\log^{m_j}(n_j+a)}{n_j+a}.
\end{equation}
These series are absolutely convergent. For a fixed outer index they
are finite nested generalized harmonic sums with logarithmic insertions.
If $r^{(d)}_{\mathbf m}=[\mathbf u^{\mathbf m}]R_d$, then
\begin{equation}\label{ord:Rcoeffrec}
 r^{(d)}_{m_1,\ldots,m_d}
 =[\mathbf u^{\mathbf m}]E_d
 +\binom{m_1+m_2+1}{m_1+1}
 r^{(d-1)}_{m_1+m_2+1,m_3,\ldots,m_d}.
\end{equation}
This is a finite-depth, exact coefficient algorithm, with base
$r^{(1)}_m=(-1)^m\gamma_m(a)/m!$.

\subsection{The diagonal regular germ and its Gamma normalization}

Write $r_d(t;a)=R_d(t,\ldots,t;a)$, with $r_0=1$.
Define the generating function, formally in $y$ and analytically for
sufficiently small $y,t$,
\begin{equation}\label{ord:Fdef}
 F_a(y,t)=\exp\left\{
 y\left(\zeta(1+t,a)-\frac1t\right)
 +\sum_{m\ge2}\frac{(-1)^{m-1}}m
                   \zeta(m+mt,a)y^m\right\}.
\end{equation}

\begin{theorem}[Diagonal regular coefficients]
\label{ord:diag}
For every $d\ge0$,
\begin{equation}\label{ord:diagformula}
 r_d(t;a)=[y^d]F_a(y,t),\qquad
 \sum_{d\ge0}r_d(0;a)y^d=\frac{\Gamma(a)}{\Gamma(a+y)}.
\end{equation}
\end{theorem}

\begin{proof}
For $\Re t>0$, the product over the common lattice gives the standard
elementary-symmetric generating series
\[
 \sum_{d\ge0}Z_d(1+t,\ldots,1+t;a)y^d
 =\prod_{n\ge0}\left(1+\frac{y}{(n+a)^{1+t}}\right)
 =\exp\sum_{m\ge1}\frac{(-1)^{m-1}}m\zeta(m+mt,a)y^m.
\]
On the diagonal, \eqref{ord:canonicalformula} becomes
$Z_d=\sum_{k=0}^dr_{d-k}(t)/(k!t^k)$.
Thus multiplication by $\exp(-y/t)$ gives \eqref{ord:Fdef},
coefficient by coefficient. Holomorphic continuation reaches $t=0$.
Finally, the Taylor series of $\log\Gamma(a+y)$ is
\[
 \log\Gamma(a)+\psi(a)y
       +\sum_{m\ge2}\frac{(-1)^m}{m}\zeta(m,a)y^m.
\]
Since $\gamma_0(a)=-\psi(a)$, exponentiation proves the last assertion.
\end{proof}

Let $B_d(a)=[y^d]\Gamma(a)/\Gamma(a+y)$.
These are explicit Bell polynomials in digamma and polygamma values.
For example, abbreviating $g=\gamma_0(a)$ and $z_m=\zeta(m,a)$,
\begin{align*}
 B_1&=g,& B_2&=\frac{g^2-z_2}{2},\\
 B_3&=\frac{g^3-3gz_2+2z_3}{6},\\
 B_4&=\frac{g^4-6g^2z_2+3z_2^2+8gz_3-6z_4}{24},\\
 B_5&=\frac{g^5-10g^3z_2+15gz_2^2+20g^2z_3
                   -20z_2z_3-30gz_4+24z_5}{120}.
\end{align*}

\subsection{All depths with one distinguished slope}

The finite part below means the coefficient of $\varepsilon^0$ after
the stated one-variable substitution. It is not a value assigned at the
intersection of polar divisors.

\begin{theorem}[A continuum of exact ordered finite parts at every depth]
\label{ord:oneslope}
Let $d\ge1$, and let $p,q\in\mathbb C$ satisfy
$p+jq\ne0$ for $0\le j\le d-1$. Put
$D_k(p,q)=\prod_{j=0}^{k-1}(p+jq)$.
Then
\begin{align}\label{ord:oneslopeformula}
 &\operatorname{FP}_{\varepsilon=0}
 Z_d(1+p\varepsilon,1+q\varepsilon,\ldots,1+q\varepsilon;a)
 \notag\\
 &\qquad=\boxed{\displaystyle B_d(a)+
 \sum_{k=1}^{d-1}\frac{q^k}{D_k(p,q)}
                      [y^{d-k}t^k]F_a(y,t).}
\end{align}
Every expression on the right is a finite polynomial in ordinary
Hurwitz zeta derivatives and generalized Stieltjes constants, with
coefficients rational in $p,q$. It uses only $\gamma_j(a)$ with
$0\le j\le d-1$ and $\zeta^{(j)}(m,a)$ with
$m\ge2$, $j\ge0$, and $m+j\le d$.
The case $q=0$ is included and gives simply $B_d(a)$.
\end{theorem}

\begin{proof}
For $k\ge1$, the tail in \eqref{ord:canonicalformula} is diagonal,
and its denominator becomes $D_k(p,q)\varepsilon^k$.
Its constant coefficient is therefore
$q^k[t^k]r_{d-k}(t;a)/D_k(p,q)$. The $k=0$ term contributes
$R_d(0)=B_d$, and the $k=d$ term is a pure pole and contributes zero.
The coefficient formula \eqref{ord:Fdef} proves the closure and the
stated finite range of orders.
\end{proof}

The same argument evaluates the complete polar part. More precisely,
for $-d\le\ell\le0$,
\begin{equation}\label{ord:allnonpositive}
 [\varepsilon^\ell]Z_d(1+p\varepsilon,1+q\varepsilon,\ldots;a)
 =\sum_{k=\max(0,-\ell)}^d
 \frac{q^{k+\ell}}{D_k(p,q)}
 [y^{d-k}t^{k+\ell}]F_a(y,t),
\end{equation}
where $D_0=1$ and $0^0=1$. For $k=0$ this formula is used only
when $\ell=0$, where its contribution is $B_d$.
Denote the finite part in \eqref{ord:oneslopeformula} by
$\mathcal F_d(p,q;a)$. For $q\ne0$ it depends only on $\lambda=p/q$:
\begin{equation}\label{ord:factorialdirection}
 \mathcal F_d(p,q;a)=B_d(a)+
 \sum_{k=1}^{d-1}\frac{[y^{d-k}t^k]F_a(y,t)}{(\lambda)_k}.
\end{equation}
Thus its dependence on the distinguished slope ratio is an explicit
finite inverse-factorial expansion. Formula \eqref{ord:factorialdirection}
can have a meromorphic continuation in $\lambda$ beyond the transverse
directions; such a continuation of a coefficient must not be confused
with restriction of the original multiple-zeta germ to a polar divisor.

This is an exact evaluation theorem, not a claim of algebraic
independence or minimality of the resulting constants. The general
ordered problem is more complicated: independent trailing slopes can
retain nonsymmetric multiple Stieltjes coefficients.

\subsection{Explicit depths two through five}

Set $g_j=\gamma_j(a)$, $g=g_0$, and
$z_m^{[j]}=\partial_s^j\zeta(s,a)|_{s=m}$; thus
$z_m=z_m^{[0]}$. Then
\begin{align}\label{ord:F2}
 \mathcal F_2&=B_2-\frac qp g_1,\\
 \mathcal F_3&=B_3-\frac qp(gg_1+z_2^{[1]})
                  +\frac{q^2g_2}{2p(p+q)}.\label{ord:F3}
\end{align}
At depth four one obtains
\begin{align}\label{ord:F4}
 \mathcal F_4={}&B_4+
 \frac qp\left(-\frac{g^2g_1}{2}+\frac{g_1z_2}{2}
                        -gz_2^{[1]}+z_3^{[1]}\right)\notag\\
 &+\frac{q^2}{p(p+q)}
       \left(\frac{g_1^2+gg_2}{2}-z_2^{[2]}\right)
 -\frac{q^3g_3}{6p(p+q)(p+2q)}.
\end{align}
At depth five,
\begin{align}\label{ord:F5}
 \mathcal F_5={}&B_5+\frac qp A_{4,1}
 +\frac{q^2}{p(p+q)}A_{3,2}
 +\frac{q^3}{p(p+q)(p+2q)}A_{2,3}
 +\frac{q^4g_4}{24p(p+q)(p+2q)(p+3q)},
\end{align}
where
\begin{align*}
 A_{4,1}={}&-\frac{g^3g_1}{6}+\frac{gg_1z_2}{2}
 -\frac{g^2z_2^{[1]}}2+\frac{z_2z_2^{[1]}}2
 -\frac{g_1z_3}{3}+gz_3^{[1]}-z_4^{[1]},\\
 A_{3,2}={}&\frac{gg_1^2}{2}+\frac{g^2g_2}{4}
 -\frac{g_2z_2}{4}+g_1z_2^{[1]}-gz_2^{[2]}
 +\frac32z_3^{[2]},\\
 A_{2,3}={}&-\frac{gg_3}{6}-\frac{g_1g_2}{2}
 -\frac23z_2^{[3]}.
\end{align*}
These formulas can be checked exactly using Newton's identities, with
independent formal symbols for every Hurwitz jet.

\subsection{Arbitrary depth-three directions and the one extra constant}

Define
\begin{equation}\label{ord:etadef}
 \eta(a)=\frac{\gamma_2(a)}2-
 \sum_{n\ge0}\frac{\gamma_1(n+a+1)+\tfrac12\log^2(n+a)}{n+a}.
\end{equation}
The summand is $O((1+\log n)/n^2)$.
The identities above imply
\[
 \eta(a)=[u]R_2(u,0;a),\qquad
 [v]R_2(0,v;a)=-\gamma_0(a)\gamma_1(a)-\zeta'(2,a)-\eta(a).
\]
For the second identity use the depth-two stuffle law, which gives
\[
 R_2(u,v;a)+R_2(v,u;a)
 =R_1(u;a)R_1(v;a)-\zeta(2+u+v,a).
\]

\begin{theorem}[Every transverse straight direction at depth three]
\label{ord:threedirections}
Let $c_1,c_2,c_3\in\mathbb C$ with
$c_1(c_1+c_2)(c_1+c_2+c_3)\ne0$. Then
\begin{align}\label{ord:threeFP}
 &\operatorname{FP}_{\varepsilon=0}
 Z_3(1+c_1\varepsilon,1+c_2\varepsilon,1+c_3\varepsilon;a)
 \notag\\
 &\quad=B_3(a)-\frac{c_3}{c_1}
       \bigl(\gamma_0(a)\gamma_1(a)+\zeta'(2,a)\bigr)
 +\frac{c_3^2\gamma_2(a)}{2c_1(c_1+c_2)}
 +\frac{c_2-c_3}{c_1}\eta(a).
\end{align}
The complete polar part is
\begin{align}\label{ord:threepoles}
 &\frac{1}{c_1(c_1+c_2)(c_1+c_2+c_3)\varepsilon^3}
 +\frac{\gamma_0(a)}{c_1(c_1+c_2)\varepsilon^2}\notag\\
 &\quad+\frac1\varepsilon
 \left\{\frac{B_2(a)}{c_1}
       -\frac{c_3\gamma_1(a)}{c_1(c_1+c_2)}\right\}.
\end{align}
\end{theorem}

\begin{proof}
At depth three the canonical identity reads
\[
 Z_3=\frac1{u_1(u_1+u_2)(u_1+u_2+u_3)}
 +\frac{R_1(u_3)}{u_1(u_1+u_2)}
 +\frac{R_2(u_2,u_3)}{u_1}+R_3(u_1,u_2,u_3).
\]
Expand $R_1$ through quadratic order and $R_2$ through linear order;
use $R_2(0,0)=B_2$, $R_3(0,0,0)=B_3$, and the two displayed
first coefficients of $R_2$. This gives both assertions.
\end{proof}

The coefficient of $\eta$ vanishes exactly on the family $c_2=c_3$
(without any assertion that $\eta$ is independent of the other constants,
or that it is nonzero at every shift).
The full labelled symmetrization also cancels it, because the two
orders of the last two slopes have opposite coefficients.
Thus generic depth-three directional finite parts need only one
nonsymmetric depth-two coefficient beyond ordinary Hurwitz data;
no new depth-three constant is required for the constant Laurent term.


% END sections/02_ordered.tex


% BEGIN sections/03_curves.tex
\section{Tangential curves and the exact amount of curve data}
\label{sec:curves}

Straight rays are only the simplest restrictions of the ordered germ.
A curve can be tangent to one or several prefix hyperplanes, producing
a higher pole order. The canonical decomposition gives a finite
calculus for all such cases, including a sharp statement about which
Taylor coefficients of the curve can affect the answer.

\subsection{A coefficient formula for every admissible analytic curve}

Let $u_i(\eps)$ be analytic near zero, with $u_i(0)=0$, and let
\begin{equation}\label{curve:prefix}
 P_j(\eps)=\sum_{i=1}^j u_i(\eps).
\end{equation}
Assume that no $P_j$ is identically zero. There are unique integers
$\nu_j\ge1$ and analytic units $p_j$ such that
\begin{equation}\label{curve:units}
 P_j(\eps)=\eps^{\nu_j}p_j(\eps),\qquad p_j(0)\ne0.
\end{equation}
Set $M_0=0$ and $M_k=\nu_1+\cdots+\nu_k$.

\begin{theorem}[Exact curve finite part and sharp jet order]
\label{curve:main}
The composed ordered function has a pole of exactly order $M_d$, with
leading coefficient $\prod_{j=1}^d p_j(0)^{-1}$. Its constant term is
\begin{equation}\label{curve:FP}
 \boxed{\displaystyle
 \FP_{\eps=0}Z_d(1+u_1(\eps),\ldots,1+u_d(\eps);a)
 =\sum_{k=0}^d[\eps^{M_k}]
 \frac{R_{d-k}(u_{k+1}(\eps),\ldots,u_d(\eps);a)}
      {p_1(\eps)\cdots p_k(\eps)}.}
\end{equation}
Regular-germ Taylor coefficients through total degree $M_d$ and curve
Taylor coefficients through order
\begin{equation}\label{curve:order}
 J=M_d+\max_{1\le j\le d}\nu_j
\end{equation}
suffice. For the class of curves with the specified prefix orders,
the bound $J$ cannot in general be decreased.
\end{theorem}

\begin{proof}
Insert \eqref{curve:units} into \eqref{ord:canonicalformula}. Its $k$th
summand is $\eps^{-M_k}$ times a holomorphic function. The constant is
therefore precisely the coefficient on the right of \eqref{curve:FP}.
Because $M_0<M_1<\cdots<M_d$, only the $k=d$ term can have pole order
$M_d$. Its numerator is $R_0=1$, so its leading coefficient is nonzero
and equals the asserted product.

In the $k$th numerator, each $u_i$ vanishes at least to first order.
Coefficients through order $M_k$ thus use only regular-germ monomials
of total degree at most $M_k$, and only curve coefficients through
that order. To invert $p_j=P_j/\eps^{\nu_j}$ through degree $M_k$
requires $P_j$ through degree $\nu_j+M_k$. These degrees are all at
most \eqref{curve:order}, proving sufficiency.

For necessity, fix an index $j$ and perturb only that prefix by
\begin{equation}\label{curve:perturb}
 P_j(\eps)\longmapsto P_j(\eps)+
                       \delta\eps^{\nu_j+M_d}.
\end{equation}
This is realized by opposite changes to $u_j,u_{j+1}$ when $j<d$,
or by a change to $u_d$ when $j=d$. All prefix orders and all other
prefixes stay fixed. The affected unit first changes in degree $M_d$.
Every denominator change in a term with $k<d$ starts after its constant
degree $M_k$; numerator changes start even later. In the $k=d$ term,
ordinary inversion gives the exact change
\begin{equation}\label{curve:sharpchange}
 \Delta\FP=-\frac{\delta}
 {p_j(0)\prod_{\ell=1}^d p_\ell(0)}.
\end{equation}
Choosing $j$ with maximal $\nu_j$ changes no curve coefficient below
order $J$ but changes the finite part by a nonzero multiple of
$\delta$. This proves sharpness.
\end{proof}

The theorem does not assign a restriction to a curve lying identically
in a polar hyperplane. Such a restriction needs an additional
operation, for example first removing a specified polar germ. A
finite algebraic value obtained by continuing a slope formula is not
by itself that extra prescription.

\subsection{Transverse curves and a depth-three fourth-jet formula}

When every prefix is transverse, $\nu_j=1$, Theorem~\ref{curve:main}
requires the curve jet through degree $d+1$. Thus a quadratic curve
description is insufficient even at depth two, and a cubic description
is insufficient in general at depth three.

An explicit example is a common reparametrization of a fixed ray.
Suppose
\[
 Z_3(1+c_1t,1+c_2t,1+c_3t;a)
 =A_{-3}t^{-3}+A_{-2}t^{-2}+A_{-1}t^{-1}+A_0+O(t),
\]
where all coefficients are given by
\eqref{ord:threeFP}--\eqref{ord:threepoles}. Put
\[
 t=h(\eps)=\eps+\alpha\eps^2+\beta\eps^3+
                         \chi\eps^4+O(\eps^5).
\]
Then exact series inversion yields
\begin{equation}\label{curve:cubicreparam}
 \boxed{\begin{aligned}
 &\FP_{\eps=0}Z_3(1+c_1h(\eps),1+c_2h(\eps),1+c_3h(\eps);a)\\
 &\quad=A_0-\alpha A_{-1}+(3\alpha^2-2\beta)A_{-2}\\
 &\qquad\quad+(-10\alpha^3+12\alpha\beta-3\chi)A_{-3}.
 \end{aligned}}
\end{equation}
For a transverse ray, $A_{-3}\ne0$, so the fourth coefficient $\chi$
has an unavoidable effect. The identity follows by taking the
constant coefficients of $h^{-1},h^{-2},h^{-3}$; their respective
values are $-\alpha$, $3\alpha^2-2\beta$, and
$-10\alpha^3+12\alpha\beta-3\chi$.

At the equal-slope ray $c_1=c_2=c_3=1$, one may insert
\[
 A_{-3}=\frac16,\qquad A_{-2}=\frac{\gamma_0(a)}2,\qquad
 A_{-1}=\frac{\gamma_0(a)^2-\gamma_1(a)-\zeta(2,a)}2,
\]
and $A_0=\mathcal F_3(1,1;a)$. This supplies a completely explicit
curved depth-three identity using ordinary Hurwitz data.

\subsection{A tangency example of arbitrarily high order}

At depth two take
\[
 u_1(\eps)=\eps,\qquad u_2(\eps)=-\eps+\eps^r,
 \qquad r\ge2.
\]
Here $(\nu_1,\nu_2)=(1,r)$, so the pole order is $r+1$ and the
sharp jet order is $2r+1$. The canonical formula simplifies exactly to
\[
 Z_2=\eps^{-r-1}+\frac{R_1(-\eps+\eps^r;a)}\eps
                  +R_2(\eps,-\eps+\eps^r;a).
\]
Consequently
\begin{equation}\label{curve:tangentexample}
 \FP Z_2=B_2(a)+\gamma_1(a).
\end{equation}
Now change $u_2$ by $\delta\eps^{2r+1}$. This leaves all lower curve
coefficients unchanged, but the first displayed term contributes an
additional constant $-\delta$. Thus the very short unperturbed answer
does not reduce the amount of curve information needed in general.

\subsection{Dependence on the ratio of straight slopes}

For $q\ne0$, put $\lambda=p/q$ and write
$\mathcal F_d(\lambda;a)=\mathcal F_d(p,q;a)$. Theorem
\ref{ord:oneslope} gives the rational function
\begin{equation}\label{curve:ratio}
 \mathcal F_d(\lambda;a)=B_d(a)+
  \sum_{k=1}^{d-1}\frac{[y^{d-k}t^k]F_a(y,t)}{(\lambda)_k}.
\end{equation}
It has at most simple poles at $0,-1,\ldots,2-d$. In particular,
for $d\ge2$,
\begin{equation}\label{curve:stieltjesresidue}
 \boxed{\displaystyle
 \operatorname*{Res}_{\lambda=2-d}\mathcal F_d(\lambda;a)
 =-\frac{\gamma_{d-1}(a)}{(d-1)!(d-2)!}.}
\end{equation}
Only the $k=d-1$ summand can contribute to this residue; its numerator
is $(-1)^{d-1}\gamma_{d-1}(a)/(d-1)!$, and the remaining factors in
the denominator have product $(-1)^{d-2}(d-2)!$.
Equation~\eqref{curve:stieltjesresidue} is an identity of the rational
slope continuation. It is not a value of the original multiple zeta
function on the excluded polar ray. It gives a useful inverse
relation: a higher Stieltjes constant can be recovered from the
slope dependence of the ordered finite-part family.

In fact the entire coefficient family has an exact inverse in terms
of these residues. For $d\ge2$, put
\[
 A_k=[y^{d-k}t^k]F_a(y,t),\qquad
 \rho_j=\operatorname*{Res}_{\lambda=-j}\mathcal F_d(\lambda;a),
 \quad 0\le j\le d-2.
\]
Partial fractions in \eqref{curve:ratio} give
\begin{equation}\label{curve:residuetransform}
 \rho_j=\frac{(-1)^j}{j!}
          \sum_{k=j+1}^{d-1}\frac{A_k}{(k-1-j)!},\qquad
 \boxed{A_k=(-1)^{k-1}\sum_{j=k-1}^{d-2}
                   \frac{j!}{(j-k+1)!}\rho_j.}
\end{equation}
To verify the inverse, set $b_j=(-1)^j j!\rho_j$. The first transform
is the upper-triangular convolution with $1/r!$; convolution with
$(-1)^r/r!$ inverts it, because $e^z e^{-z}=1$. This proves the second
formula and recovers \eqref{curve:stieltjesresidue} at $k=d-1$.

\subsection{Recovering every regular jet from directional finite parts}

The one-distinguished-slope family determines diagonal combinations
of regular coefficients. Allowing all independent slopes recovers
each coefficient that can appear, with no loss of information.

Write $H_{r,m}(\boldsymbol v;a)=[t^m]R_r(t\boldsymbol v;a)$ for the
homogeneous degree-$m$ Taylor polynomial, and let
\[
 \mathcal G_d(\boldsymbol c;a)=\FP_{\eps=0}
 Z_d(1+c_1\eps,\ldots,1+c_d\eps;a).
\]
For $P_j=c_1+\cdots+c_j\ne0$, the canonical formula gives
\begin{equation}\label{curve:generalray}
 \mathcal G_d(\boldsymbol c;a)=B_d(a)+
 \sum_{k=1}^{d-1}
 \frac{H_{d-k,k}(c_{k+1},\ldots,c_d;a)}{P_1\cdots P_k}.
\end{equation}

\begin{corollary}[Exact reconstruction from directional data]
\label{curve:reconstruction}
The rational function $\mathcal G_d$ uniquely determines every
$H_{d-k,k}$, $1\le k\le d-1$, and its constant $B_d$. In particular,
the full collection of directional finite parts at all depths
determines every regular germ $R_r$: its homogeneous degree-$m$ part
is recovered at depth $r+m$.
\end{corollary}

\begin{proof}
Use the independent prefix coordinates $P_1,\ldots,P_d$ and descend
from $k=d-1$ to $k=1$. Suppose all higher-index summands in
\eqref{curve:generalray} have been recovered and subtracted. Take the
coefficient of $P_1^{-1}\cdots P_k^{-1}$ in the remainder, treating
$P_{k+1},\ldots,P_d$ as generic parameters. The lower-index summands
are holomorphic in at least one of these first $k$ variables and
contribute zero. The surviving coefficient is
\[
 H_{d-k,k}(P_{k+1},P_{k+2}-P_{k+1},\ldots,P_d-P_{d-1};a).
\]
This invertible change of the remaining variables recovers the
entire polynomial $H_{d-k,k}$. Subtracting all recovered terms leaves
$B_d$. For a general $R_r$ and $m\ge1$, choose $d=r+m$ and $k=m$.
The degree-zero part is $R_r(0)=B_r$, already known from the Gamma
quotient. No higher Taylor degree than those in
\eqref{curve:generalray} can affect a finite part at the fixed depth.
\end{proof}

% END sections/03_curves.tex


% BEGIN sections/04_harmonic_polylog.tex
\section{The harmonic coefficient and its polylogarithmic representations}
\label{sec:harmonic}

The coefficient $\eta(a)$ in \eqref{ord:etadef} is not an undefined
remainder. It has an absolutely convergent Stieltjes series, belongs
to an established harmonic-zeta convention, and admits a convergent
Mellin representation with an elementary kernel when $a=1$.
This section also records exact difference and derivative identities.

\subsection{Strict and inclusive harmonic conventions}

For $n\ge0$ define the strict shifted harmonic sum
\[
 H_n^{(s)}(a)=\sum_{k=0}^{n-1}(k+a)^{-s},\qquad H_0^{(s)}(a)=0.
\]
It is entire in $s$ for fixed $n$. Its spectral derivatives are
finite logarithm-weighted harmonic sums. Summing first over the
inner index gives
\begin{equation}\label{harm:strict}
 Z_2(s,1;a)=\sum_{n\ge0}\frac{H_n^{(1)}(a)}{(n+a)^s}.
\end{equation}
The canonical decomposition with the second deviation zero yields
\begin{equation}\label{harm:etaLaurent}
 Z_2(1+u,1;a)=\frac1{u^2}+\frac{\gamma_0(a)}u+
 B_2(a)+\eta(a)u+O(u^2).
\end{equation}
Indeed the regular part is $R_2(u,0;a)$ and its first coefficient is
exactly \eqref{ord:etadef}.

Karg\i n et al. use the inclusive harmonic Hurwitz function
\cite{Kargin}
\[
 \zeta_H(s,a)=\sum_{n\ge0}(n+a)^{-s}
                           \sum_{k=0}^{n}(k+a)^{-1}.
\]
The relation and the coefficient conversion are
\begin{align}
 \zeta_H(s,a)&=Z_2(s,1;a)+\zeta(s+1,a),\label{harm:inclusive}\\
 \gamma_H(0,a)&=\frac{\gamma_0(a)^2+\zeta(2,a)}2,\notag\\
 \eta(a)&=\boxed{-\gamma_H(1,a)-\zeta'(2,a).}\label{harm:conversion}
\end{align}
Here $\gamma_H(m,a)$ is normalized by the coefficient
$(-1)^m\gamma_H(m,a)/m!$ in the regular Laurent series of
$\zeta_H(1+u,a)$. This fixes both the strict/inclusive difference and
the sign of the first coefficient.

\subsection{Shift and parameter-derivative identities}

\begin{proposition}[Transport of the ordering coefficient]
\label{harm:transport}
For every $a>0$,
\begin{equation}\label{harm:shift}
 \boxed{\eta(a+1)-\eta(a)=\frac{\gamma_1(a+1)}a.}
\end{equation}
Writing $\partial_1,\partial_2$ for the spectral derivatives of
$Z_2$, one also has the convergent identity
\begin{align}\label{harm:derivative}
 \eta'(a)={}&-Z_2(2,1;a)-\partial_1Z_2(2,1;a)
                 +\partial_2Z_2(2,1;a)\notag\\
 &+\gamma_1(a)\zeta(2,a)+\zeta'(3,a).
\end{align}
All three double sums on the right converge absolutely, including
their indicated logarithmic insertions.
\end{proposition}

\begin{proof}
Removing the terms with smallest index zero gives
\[
 Z_2(s,1;a)-Z_2(s,1;a+1)=a^{-1}\zeta(s,a+1).
\]
The coefficient of $s-1$ is
$\eta(a)-\eta(a+1)=-\gamma_1(a+1)/a$, proving
\eqref{harm:shift}.

In an absolute-convergence domain differentiation in the common
shift gives
\[
 \partial_a Z_2(s,1;a)=-sZ_2(s+1,1;a)-Z_2(s,2;a).
\]
Use the stuffle identity
\[
 Z_2(s,2;a)=\zeta(s,a)\zeta(2,a)-Z_2(2,s;a)-\zeta(s+2,a)
\]
and continue near $s=1$. Taking the coefficient of $s-1$ in
\eqref{harm:etaLaurent} gives \eqref{harm:derivative}.
Convergence follows from $H_n^{(1)}(a)=O(1+\log n)$ and the extra
outer square; spectral derivatives add only logarithmic factors.
\end{proof}

For integer shifts, repeated use of \eqref{harm:shift} yields the
finite identity
\[
 \eta(a+N)=\eta(a)+\sum_{k=0}^{N-1}
                         \frac{\gamma_1(a+k+1)}{a+k}.
\]
These transport identities are consistent with the harmonic Hurwitz
shift recurrence; they are included to make the ordering coefficient
usable without a change of convention.

\subsection{A convergent Lerch--polylogarithm Mellin formula}

For $0<z<1$, let $\Phi(z,s,a)=\sum_{n\ge0}z^n/(n+a)^s$.
The harmonic generating function is
\begin{equation}\label{harm:gen}
 \sum_{n\ge0}H_n^{(1)}(a)z^n=\frac{z\Phi(z,1,a)}{1-z}.
\end{equation}
It follows by interchanging two absolutely convergent sums. Hence
\begin{equation}\label{harm:Mellin}
 \Gamma(s)Z_2(s,1;a)=\int_0^\infty t^{s-1}K_a(t)\,dt,
 \quad K_a(t)=\frac{e^{-(a+1)t}\Phi(e^{-t},1,a)}{1-e^{-t}},
 \quad \Re s>1.
\end{equation}

Put $h_a=\gamma_0(a)-\gamma$ and define
\begin{align}\label{harm:J}
 J_a(s)={}&\int_0^1t^{s-1}
       \left(K_a(t)-\frac{-\log t+h_a}{t}\right)\,dt\notag\\
 &+\int_1^\infty t^{s-1}K_a(t)\,dt.
\end{align}
The kernel subtraction is $O(1+|\log t|)$ at zero. Indeed, put
$d_n=(n+a)^{-1}-(n+1)^{-1}$. Then $d_n=O((n+1)^{-2})$,
$\sum d_n=h_a$, and
\[
 \sum_{n\ge0}\frac{|1-e^{-nt}|}{(n+1)^2}
        =O\bigl(t(1+|\log t|)\bigr).
\]
Comparing $\Phi(e^{-t},1,a)$ with $\Phi(e^{-t},1,1)$ proves the
asserted estimate, locally uniformly for $a>0$.
Thus $J_a$ is holomorphic for
$\Re s>0$, with all fixed spectral derivatives obtained by inserting
powers of $\log t$. We have the exact continuation
\begin{equation}\label{harm:Jidentity}
 \Gamma(s)Z_2(s,1;a)=\frac1{(s-1)^2}+
                 \frac{h_a}{s-1}+J_a(s).
\end{equation}

\begin{proposition}[A normally convergent integral for $\eta$]
\label{harm:etaIntegral}
Let $g_0=\gamma_0(a)$. Then
\begin{align}\label{harm:etaintegral}
 \eta(a)={}&J_a'(1)+\gamma B_2(a)
 -\frac{\gamma^2+\zeta(2)}2g_0
 +\frac{\gamma^3}6+\frac{\gamma\zeta(2)}2+\frac{\zeta(3)}3,\\
 J_a(1)={}&\frac{(g_0-\gamma)^2+\zeta(2)-\zeta(2,a)}2.\label{harm:Jone}
\end{align}
\end{proposition}
\begin{proof}
Multiply \eqref{harm:etaLaurent} by the Taylor expansion of
$\Gamma(1+u)$ through cubic order and compare the constant and linear
coefficients with \eqref{harm:Jidentity}. The expansion of
$\log\Gamma(1+u)$ gives
\[
 [u^2]\Gamma(1+u)=\frac{\gamma^2+\zeta(2)}2,\quad
 [u^3]\Gamma(1+u)=-\frac{\gamma^3}6-
                  \frac{\gamma\zeta(2)}2-\frac{\zeta(3)}3,
\]
which proves the two formulas.
\end{proof}

At $a=1$, $\Phi(z,1,1)=\Li_1(z)/z=-\log(1-z)/z$ and
\[
 K_1(t)=-\frac{\log(1-e^{-t})}{e^t-1},\qquad J_1(1)=0.
\]
The integral formula reduces to
\begin{align}\label{harm:etaone}
 \eta(1)={}&\frac{\gamma^3}6-\frac{\gamma\zeta(2)}2+
                \frac{\zeta(3)}3\notag\\
 &+\int_0^1\log t\left(
      -\frac{\log(1-e^{-t})}{e^t-1}+\frac{\log t}{t}\right)dt
 +\int_1^\infty\frac{-\log t\log(1-e^{-t})}{e^t-1}\,dt.
\end{align}
Both integrals converge ordinarily. Numerically,
$\eta(1)=0.5367870243584185675036078843488117\ldots$.
The numerical section compares this integral with Cauchy extraction
from an independently continued ordered double sum.

All higher strict harmonic coefficients follow from the same kernel;
they do not require a separate continuation. Put
\[
 e_j=[u^j]\frac1{\Gamma(1+u)},\qquad
 C_m(a)=[u^m]R_2(u,0;a).
\]
Multiplying \eqref{harm:Jidentity} by $1/\Gamma(1+u)$ gives, for
every integer $m\ge0$,
\begin{align}\label{harm:allcoefficients}
 C_m(a)&=e_{m+2}+h_a e_{m+1}
       +\sum_{j=0}^m\frac{e_{m-j}}{j!}J_a^{(j)}(1),\\
 C_m(a)&=\frac{(-1)^m\gamma_H(m,a)-\zeta^{(m)}(2,a)}{m!}.
                  \label{harm:allconversion}
\end{align}
The second formula is \eqref{harm:inclusive} coefficient by
coefficient. Each $J_a^{(j)}(1)$ is an ordinarily convergent integral
obtained by inserting $\log^j t$ in \eqref{harm:J}. Thus the complete
strict harmonic Stieltjes hierarchy has a single convergent kernel,
with reciprocal-Gamma coefficients supplying its finite normalization.

\subsection{Polylogarithmic primitives and harmonic coefficients}

For $|z|<1$, the strict multiple polylogarithm
\[
 \Li_{s_1,\ldots,s_d}(z)=
 \sum_{n_1>\cdots>n_d\ge1}
       \frac{z^{n_1}}{n_1^{s_1}\cdots n_d^{s_d}}
\]
has the ordinary normalized Euler primitive
\begin{equation}\label{harm:primitive}
 \int_0^z\Li_{s_1,s_2,\ldots,s_d}(w)\,\frac{dw}{w}
       =\Li_{s_1+1,s_2,\ldots,s_d}(z).
\end{equation}
All fixed spectral derivatives commute with this identity on compact
subdisks. Their inner coefficients are the unshifted logarithm-weighted
harmonic sums that also occur in the remainder construction
\eqref{ord:Ecoeff}. At all indices one,
\begin{equation}\label{harm:allones}
 \Li_{\{1\}^d}(z)=\frac{[-\log(1-z)]^d}{d!},\qquad
 \sum_{d\ge1}y^{d-1}\Li_{s,\{1\}^{d-1}}(z)
 =\sum_{n\ge1}\frac{z^n}{n^s}
             \frac{\Gamma(n+y)}{\Gamma(n)\Gamma(1+y)}.
\end{equation}
The second identity follows by taking elementary symmetric
polynomials of $1,1/2,\ldots,1/(n-1)$; the Gamma quotient is the
polynomial $\prod_{j=1}^{n-1}(1+y/j)$, including at its apparently
exceptional parameters. One may first take $y$ near zero and then
continue in $y$. These are classical
polylogarithmic and harmonic generating identities, recorded here
to give a concrete derivative and antiderivative dictionary for the
Laurent coefficients. The extension at the singular endpoint is the
specified spectral calculus proved above; the endpoint is not obtained
by inserting $z=1$ into a divergent series.

% END sections/04_harmonic_polylog.tex


% BEGIN sections/05_cubic.tex
\section{Cubic coincident Stieltjes products and a diagonal Tornheim jet}
\label{sec:cubic}

The incoming \emph{Coincident Stieltjes} article explicitly asks for the
unit-coordinate finite part of $\int_0^1\psi(x)^3\,dx$ and an all-index
three-factor closure theorem.  It suggests higher-depth coordinates and
does not assume a reduction to ordinary zeta constants.  The results
below supply an exact reduction to a specified diagonal Tornheim
derivative, together with a holomorphic germ for all mixed Stieltjes
indices.  They do not evaluate that remaining derivative in a
depth-one constant algebra.

The triple-Hurwitz/Tornheim correspondence is classical, as is the
Mellin--polylogarithm continuation used to compute Tornheim functions.
Relevant antecedents include Espinosa--Moll, Borwein--Dilcher,
Bailey--Borwein, and Onodera\cite{EMTornheim,BorweinDilcher,BaileyBorwein,Onodera2021}.
The low diagonal jets below are recovered identities from this established
setting.

\subsection{Normalization and the classical three-factor kernel}

Use the unit coordinate $x$ and, for $e\in\mathbb Z$ and
$k\in\mathbb Z_{\ge0}$, the convention
\[
 \operatorname{FP}\int_0^1 x^e\log^k x\,dx=
 \begin{cases}(-1)^k k!/(e+1)^{k+1},&e\ne-1,\\0,&e=-1.
 \end{cases}
\]
This agrees with the constant term after removing the divergent
endpoint powers and logarithms from $\int_\epsilon^1$.
Let
\[
 g(u,x)=\zeta(1+u,x)-\frac1u
       =\sum_{m\ge0}\frac{(-1)^m}{m!}\gamma_m(x)u^m,
 \qquad
 b(u)=\zeta(1+u)-\frac1u,
 \qquad
 a(u)=-(1+u)\zeta(2+u).
\]
The removable singularities in $g$ and $b$ are always removed.
In particular $g(0,x)=-\psi(x)$, $b(0)=\gamma$, and
$a(0)=-\zeta(2)$.

Initially in its absolutely convergent region, put
\[
 T(A,B,C)=\sum_{m,n\ge1}\frac1{m^A n^B(m+n)^C}.
\]
Its meromorphic continuation is understood in every later formula.
The initial absolute-convergence domain is given by
$\Re(A+C)>1$, $\Re(B+C)>1$, and $\Re(A+B+C)>2$.
Write $K_2(s,t)=\int_0^1\zeta(s,x)\zeta(t,x)\,dx$ and
$K_3(s_1,s_2,s_3)=\int_0^1\prod_{i=1}^3\zeta(s_i,x)\,dx$ in an
initial convergence domain, then continue meromorphically.
If $\lambda_i=1-s_i$, the classical triple kernel is
\begin{align}
 K_3(s_1,s_2,s_3)
 &=\frac{2\Gamma(\lambda_1)\Gamma(\lambda_2)\Gamma(\lambda_3)}
          {(2\pi)^{\lambda_1+\lambda_2+\lambda_3}}
 \sum_{k=1}^3
 \cos\frac{\pi(\lambda_i+\lambda_j-\lambda_k)}2\,
 T(\lambda_i,\lambda_j,\lambda_k),             \tag{C1}\label{cub:C1}
\end{align}
where $i<j$ are the indices other than $k$.
For $\Re\lambda_i>1$, the Fourier series
\[
 \zeta(1-\lambda,x)=\Gamma(\lambda)(2\pi)^{-\lambda}
 \sum_{n\ne0}|n|^{-\lambda}
       e^{-\pi i\lambda\operatorname{sgn}(n)/2}e^{2\pi inx}
\]
converges absolutely.  The zero-frequency condition in the product
requires two positive frequencies and one negative, or its opposite.
The two sign regions paired for a fixed negative position give the
cosine in (C1). This proves (C1) in a nonempty open set; meromorphic
continuation gives the stated identity. The same argument gives the pair kernel
\[
K_2(s,t)=\frac{2\Gamma(1-s)\Gamma(1-t)}{(2\pi)^{2-s-t}}
\cos\frac{\pi(s-t)}2\,\zeta(2-s-t).
\]

\subsection{The complete cubic subtraction}

For compactness write $K_{ij}=K_2(1+u_i,1+u_j)$,
$K_{123}=K_3(1+u_1,1+u_2,1+u_3)$, and $U=u_1+u_2+u_3$.
Define the meromorphic germ
\[
 I_3(\boldsymbol u)=K_{123}-\sum_{k=1}^3\frac{K_{ij}}{u_k}
                    -\frac1{u_1u_2u_3}.        \tag{C2}\label{cub:C2}
\]
The mean of a single Hurwitz factor is zero in the initial domain, so
this is precisely the meromorphic continuation of
$\int_0^1\prod_i g(u_i,x)\,dx$.

\begin{theorem}[Cubic holomorphic completion]
The germ
\[
 F_3(\boldsymbol u)=I_3(\boldsymbol u)
   +\sum_{k=1}^3\frac{a(u_k)}{u_i+u_j}
   +\sum_{k=1}^3\frac{b(u_i)b(u_j)}{u_k}       \tag{C3}\label{cub:C3}
\]
is holomorphic in a neighborhood of $\boldsymbol u=0$. For all
$m_1,m_2,m_3\ge0$,
\[
 \operatorname{FP}\int_0^1\prod_{i=1}^3\gamma_{m_i}(x)\,dx
 =(-1)^{m_1+m_2+m_3}\prod_{i=1}^3m_i!\,
   [u_1^{m_1}u_2^{m_2}u_3^{m_3}]F_3(\boldsymbol u). \tag{C4}\label{cub:C4}
\]
\end{theorem}

\begin{proof}
The shift relation gives
\[
 g(u,x)=x^{-1-u}+h(u,x),\qquad h(u,x)=g(u,1+x),
 \qquad h(u,x)=b(u)+a(u)x+O(x^2).
\]
All these Taylor statements are uniform on small compact spectral
neighborhoods.  Write $h_i=h(u_i,x)$.  The product of the three singular
terms contributes $-1/(2+U)$ after meromorphic integration. A term with
the two singular factors in positions $i,j$ and smooth factor $h_k$
contributes
\[
 -\frac{b(u_k)}{1+u_i+u_j}-\frac{a(u_k)}{u_i+u_j}
 +\int_0^1x^{-2-u_i-u_j}
             \bigl(h_k-b(u_k)-a(u_k)x\bigr)\,dx.
\]
A term with just the singular factor in position $k$ contributes
\[
 -\frac{b(u_i)b(u_j)}{u_k}
 +\int_0^1x^{-1-u_k}\bigl(h_i h_j-b(u_i)b(u_j)\bigr)\,dx.
\]
The remaining term is $\int_0^1h_1h_2h_3\,dx$.
After (C3), all remaining integrals are normally convergent near zero:
their endpoint majorants have the form
$Cx^{-\eta}(1+|\log x|^M)$, with $\eta<1$, after any prescribed finite
number of spectral derivatives.  This proves holomorphy.

The deleted expressions are exactly the integrations of resonant
monomials with exponent $-1$ at $\boldsymbol u=0$. Each coefficient of
such a monomial is a multiple of $x^{-1}\log^j x$, whose unit-coordinate
finite part is zero. The other integrated monomials have precisely the
finite parts prescribed above. Coefficient extraction under the
normally convergent remainder proves (C4).
\end{proof}

The mixed spectral derivatives in (C4) apply to the completed germ.
Uncompleted individual Tornheim summands can have intersecting polar
hyperplanes at the origin. In particular one must not read (C4) as a
formula involving unspecified ordinary partial derivatives of
$T$ at $(0,0,0)$.

\subsection{Every argument derivative}

The same proof gives a finite all-order extension. For
$p_i\in\mathbb Z_{\ge0}$, let
\begin{gather*}
 c_p(u)=(-1)^p(1+u)_p,\qquad
 h_i(u_i,x)=\partial_x^{p_i}g(u_i,1+x),\\
 G_i=\partial_x^{p_i}g(u_i,x)
     =c_{p_i}(u_i)x^{-p_i-1-u_i}+h_i(u_i,x).
\end{gather*}
Let $I_{\boldsymbol p}$ be the meromorphic product integral of the
$G_i$. With $s_i=1+p_i+u_i$, it is explicitly
\begin{align*}
 I_{\boldsymbol p}={}&\left(\prod_{i=1}^3 c_{p_i}(u_i)\right)
                         K_3(s_1,s_2,s_3)\\
 &-\sum_{k=1}^3\frac{\delta_{p_k0}c_{p_i}(u_i)c_{p_j}(u_j)}{u_k}
                         K_2(s_i,s_j)
 -\frac{\prod_{i=1}^3\delta_{p_i0}}{u_1u_2u_3}.
\end{align*}
The single-Hurwitz integrals vanish in the initial convergence domain
and hence identically after continuation.
Then the holomorphic completion is
\begin{align}
 F_{\boldsymbol p}=I_{\boldsymbol p}
 &+\sum_{k=1}^3\frac{c_{p_k}(u_k)}{u_k}
       [x^{p_k}]\bigl(h_i(u_i,x)h_j(u_j,x)\bigr)\notag\\
 &+\sum_{k=1}^3\frac{c_{p_i}(u_i)c_{p_j}(u_j)}{u_i+u_j}
       [x^{p_i+p_j+1}]h_k(u_k,x).             \tag{C5}\label{cub:C5}
\end{align}
Its coefficients generate
$\operatorname{FP}\int_0^1\prod_i\gamma_{m_i}^{(p_i)}(x)\,dx$ with
the same factorial and sign in (C4). Indeed a subset $S$ of singular
factors has exponent $-|S|-\sum_{i\in S}p_i-\sum_{i\in S}u_i$;
the unique resonant smooth Taylor degree is
$|S|+\sum_{i\in S}p_i-1$.  For a singleton or pair this gives precisely
(C5). For all three factors the needed smooth degree is positive but
there is no smooth factor, so no further resonant term occurs.
To justify the normal convergence rather than merely identify the
formal poles, subtract every smooth Taylor term up to that threshold.
The integrated lower Taylor degrees have nonzero integer denominators
at the origin, so they are holomorphic and retain their prescribed
monomial finite parts. The only denominators vanishing at the origin
are exactly those displayed in (C5). The remaining Taylor remainders
are integrable uniformly on sufficiently small compact spectral
neighborhoods, also after any fixed finite number of derivatives.
Every smooth coefficient is a finite combination of Hurwitz jets,
because
\[
 [x^\ell]h_i(u_i,x)
 =\frac{c_{p_i+\ell}(u_i)}{\ell!}
      \zeta(1+p_i+\ell+u_i)
   -\frac{\delta_{p_i0}\delta_{\ell0}}{u_i}.
\]
The Pochhammer coefficient expansions are the usual finite harmonic
polynomials. Thus (C5) is an explicit finite recipe for all cubic
argument-derivative moments.

\subsection{An independently defined diagonal Tornheim function}

Let $\tau(s)=T(s,s,s)$ initially for $\Re s>2/3$, and then continue
this one-variable function.  The diagonal restriction is essential:
$T$ itself has no direction-independent value at the multivariate
origin. For example, $T(s,s,0)=\zeta(s)^2$ tends to $1/4$, whereas
the diagonal limit is $1/3$, as recovered in (C8) below.

Here is a convergent formula defining the germ of $\tau$ near zero
without using any Hurwitz product integral. Set
$L_s(t)=\operatorname{Li}_s(e^{-t})$, $G(s)=\Gamma(1-s)$, and
\begin{align*}
 R_s(t)={}&L_s(t)^2-G(s)^2t^{2s-2}
   -2G(s)\zeta(s)t^{s-1}
   +2G(s)\zeta(s-1)t^s-\zeta(s)^2,\\
 E(s)={}&\int_0^1t^{s-1}R_s(t)\,dt
          +\int_1^\infty t^{s-1}L_s(t)^2\,dt.
\end{align*}
For sufficiently small $\delta>0$, $E$ is holomorphic when $|s|<\delta$.
Indeed the Jonqui\`ere expansion \cite{DLMFPolylog}
\[
 L_s(t)=\Gamma(1-s)t^{s-1}
        +\sum_{k=0}^\infty\frac{(-1)^k}{k!}\zeta(s-k)t^k,
 \qquad |t|<2\pi,
\]
shows that the subtracted integrand is bounded by
$C(t^{-2\delta}+t^{-\delta})(1+|\log t|^M)$ after any fixed number
of spectral derivatives. The second integral converges exponentially.
Direct Mellin integration first for $2/3<\Re s<1$ and then continuation
gives the holomorphic formula
\[
 \tau(s)=\frac1{\Gamma(1+s)}\left\{
 sE(s)+\frac{s\Gamma(1-s)^2}{3s-2}
 +\frac{2s\Gamma(1-s)\zeta(s)}{2s-1}
 -\Gamma(1-s)\zeta(s-1)+\zeta(s)^2\right\}. \tag{C6}\label{cub:C6}
\]
This is the $\theta=1$ local-subtraction form of the classical
Crandall Mellin continuation. In particular $\tau'''(0)$ is an
unambiguous, independently computable spectral constant.

\subsection{The cubic digamma reduction}

Put $\ell=\log(2\pi)$, and define
\[
 Z(t)=t\zeta(1+t),\quad
 A_d(t)=\frac{\Gamma(1-t)^2}{\Gamma(1-2t)\cos(\pi t)},\quad
 B(t)=\Gamma(1-t)^3(2\pi)^{3t}\cos(\pi t/2).
\]
The classical pair formula yields
$K_2(1+t,1+t)=-A_d(t)Z(2t)/t^2$, while (C1) gives
$K_3(1+t,1+t,1+t)=-6B(t)\tau(-t)/t^3$. Hence
\[
 F_3(t,t,t)=
 \frac{-6B(t)\tau(-t)+3A_d(t)Z(2t)-1+3(Z(t)-1)^2}{t^3}
 +\frac{3a(t)}{2t}.                           \tag{C7}\label{cub:C7}
\]
Since the left side is holomorphic, comparison of the three negative
powers gives
\[
 \tau(0)=\frac13,\qquad
 \tau'(0)=\ell,\qquad
 \tau''(0)=3\ell^2-2\gamma^2-4\gamma_1+\zeta(2).\tag{C8}\label{cub:C8}
\]
These are recovered low-jet identities; the first two are well known,
and Onodera's 2021 paper evaluates second derivatives at the origin.

For completeness the needed expansions are
\begin{align*}
 Z(t)&=1+\gamma t-\gamma_1t^2+\tfrac12\gamma_2t^3+O(t^4),\\
 A_d(t)&=1+2\zeta(2)t^2-2\zeta(3)t^3+O(t^4),\\
 \log B(t)&=3(\gamma+\ell)t+\tfrac34\zeta(2)t^2
                 +\zeta(3)t^3+O(t^4),\\
 a'(0)&=-\zeta(2)-\zeta'(2).
\end{align*}
Taking the constant coefficient of (C7) and using
$F_3(0,0,0)=-\operatorname{FP}\int_0^1\psi(x)^3\,dx$ proves
\begin{align}
 \operatorname{FP}\int_0^1\psi(x)^3\,dx={}&-\tau'''(0)
 -9\gamma^3-18\gamma^2\ell-30\gamma\gamma_1
 -36\ell\gamma_1-12\gamma_2+9\ell^3\notag\\
 &+\frac32\gamma\zeta(2)+9\ell\zeta(2)
 +\frac32\zeta(2)+8\zeta(3)+\frac32\zeta'(2). \tag{C9}\label{cub:C9}
\end{align}
An equivalent unsimplified coefficient certificate is
\begin{align*}
 F_3(0,0,0)={}&-6[t^3]B(t)\tau(-t)
 +12\gamma_2+12\gamma\zeta(2)-6\zeta(3)-6\gamma\gamma_1\\
 &-\tfrac32\bigl(\zeta(2)+\zeta'(2)\bigr).
\end{align*}

\subsection{Directional Laurent constants do not fix the coordinate finite part}

Let $\alpha_1,\alpha_2,\alpha_3$ be fixed complex numbers with
$\alpha_i\ne0$ and $\alpha_i+\alpha_j\ne0$. If $\operatorname{CT}$
denotes the one-variable Laurent constant and
$Q=\operatorname{FP}\int_0^1\gamma_0(x)^3\,dx$, then (C3) gives
\begin{align}
 Q={}&\operatorname{CT}_{t=0} I_3(\alpha_1t,\alpha_2t,\alpha_3t)
 -\bigl(\zeta(2)+\zeta'(2)\bigr)
       \sum_{k=1}^3\frac{\alpha_k}{\alpha_i+\alpha_j}\notag\\
 &-\gamma\gamma_1\sum_{k=1}^3\frac{\alpha_i+\alpha_j}{\alpha_k}.
                                                        \tag{C10}\label{cub:C10}
\end{align}
For the diagonal this correction is
$-\frac32(\zeta(2)+\zeta'(2))-6\gamma\gamma_1$.
This is a normalization warning, not an identified false claim in an
inspected source. Removing only the negative Laurent powers after
diagonal specialization would omit the finite coefficients of the
resonant expressions. The full germ (C3) removes exactly those entire
expressions and yields the fixed-coordinate finite part.

\subsection{Independent numerical diagnostic}

The script \texttt{verify\_cubic.py} computes the first four diagonal
Tornheim derivatives from (C6), using the convergent Jonqui\`ere series
on $(0,1)$ and an exponentially convergent integral on $(1,\infty)$.
The independent digamma calculation integrates
\[
 \psi(x)^3+x^{-3}+3\gamma x^{-2}
                -3\bigl(\zeta(2)-\gamma^2\bigr)x^{-1}
\]
ordinarily and then adds $1/2+3\gamma$. Near zero it evaluates the
regular remainder by the convergent Taylor series for $\psi(1+x)$,
avoiding cancellation of huge singular terms.
At 35 working decimal digits and 50 Jonqui\`ere terms the values are
\begin{align*}
 \tau'''(0)&\approx86.657866601100073652903276858762183,\\
 \operatorname{FP}\int_0^1\psi(x)^3\,dx
 &\approx1.9662627343915343406643753384116314.
\end{align*}
The two independent expressions in (C9) differ by less than
$4.1\times10^{-34}$. This is a floating-point diagnostic, not an
interval enclosure or substitute for the proof above.


% END sections/05_cubic.tex


% BEGIN sections/06_verification.tex
\section{Verification, proposed status updates, and integration}
\label{sec:verification}

\subsection{What the proofs establish and what the computations check}

The analytic arguments establish the normal convergence of the tail
remainders, the meromorphic identities, the exact curve-jet bound,
and the endpoint completion at every stated index. These conclusions
do not depend on numerical precision. The attached verification has
two distinct roles: exact algebra checks the printed coefficients,
and independent numerical constructions test signs, normalizations,
and representative specializations.

The primary exact calculation uses Newton's elementary-symmetric
recurrence with a separate formal symbol for every Hurwitz jet. It
does not obtain the displayed coefficients by differentiating a copy
of the displayed exponential. Sixty checks cover the Gamma polynomials
$B_1,\ldots,B_5$, the explicit formulas through depth five, the
equal-slope Laurent coefficients through depth seven, the zero-trailing
slope family, and the complete six-permutation depth-three stuffle
identity. A separately written calculation performs seven additional
ordered-coefficient checks. Five further exact checks verify the
cubic pole cancellations, the complete constant in (C9), and the
Mellin counterterm exponents. The seventy-two checks all pass.

For the ordered numerical diagnostic, the double and triple sums are
continued directly by summing a finite initial range and attaching
Euler--Maclaurin tails of the ordinary Hurwitz factors. Cauchy
coefficient extraction then obtains the finite part on a small
complex circle. This calculation does not use the regular-germ
recursion to evaluate the left side. At $a=1$, the additional
coefficient $\eta$ is also computed from the independently convergent
integral \eqref{harm:etaone}. The cubic-product diagnostic separately
uses the polylogarithmic Mellin formula and the spatially subtracted
digamma integral described in Section~\ref{sec:cubic}.

\begin{center}
\small
\begin{tabular}{p{.44\textwidth}p{.42\textwidth}}
\toprule
Diagnostic & Representative value\\
\midrule
$\eta(1)$ & $0.536787024358418567503607884349\ldots$\\[3pt]
$\mathcal F_3(1,\tfrac12;1)$ &
 $0.446979185190867817804303675619\ldots$\\[3pt]
$\FP Z_3(1+\tfrac54\eps,1+\tfrac12\eps,
1+\tfrac34\eps;\tfrac34)$ &
 $0.169462871179343781164830292059\ldots$\\[3pt]
$\tau'''(0)$ & $86.6578666011000736529032768588\ldots$\\[3pt]
$\FP\int_0^1\psi(x)^3\,dx$ &
 $1.96626273439153434066437533841\ldots$\\
\bottomrule
\end{tabular}
\end{center}

The recorded ordered run uses 50 working decimal digits, twenty
Cauchy nodes, a finite cutoff of 36, and twenty Euler--Maclaurin terms.
The direct triple evaluation uses forty additional guard digits to
protect the small high-order Hurwitz tails before their multiplication
by large Euler--Maclaurin coefficients.
Across the four depth-three comparisons at $a=1$ and $a=3/4$, the
largest absolute residual is below $8\times10^{-35}$; the independent
$\eta(1)$ comparison is below $10^{-38}$. The digamma-cube comparison
uses 35 working digits and fifty Jonqui\`ere terms, with residual
below $4.1\times10^{-34}$. The saved JSON records the actual complex
values and residuals. These are floating-point diagnostics, with
truncation and roundoff; the number of working digits is not a claim
that every printed digit is certified. No integer-relation search is
used as proof, and no interval enclosure is asserted.

\subsection{Independent review and limits of the source audit}

The ordered and cubic derivations were reviewed independently of
their initial derivation, including the tail majorants, divided
differences, factorials, Fourier phases, endpoint resonances, and the
third-order Tornheim expansion. A separate review checked the sharp
curve perturbation, the slope residue, the harmonic conversion, and
the Mellin kernel. Review led to a more explicit argument-derivative
kernel: both the triple and pair kernels carry their shifted spectral
arguments. It also led to an explicit remainder subtraction through
the full integrability threshold. Those clarifications are incorporated
in the proofs, rather than left as qualifications to the formulas.

No new concrete false mathematical claim was identified in the
inspected canonical or incoming material. In particular, the inspected
quadratic coincident completion and the recent Dougall resonance
formulas were consistent with their stated scope. This is a targeted
audit, not a claim that the hundreds of pages in the canonical
manuscript have all been re-proved. Several tempting proposed
extensions were already present in the incoming reports and have not
been relabelled as new results here.

The changes recommended for the repository are therefore precise
status and normalization updates:
\begin{enumerate}[leftmargin=*,itemsep=4pt]
\item Mark the independent-exponent ordered depth-three question in
the resonant-jets report as answered by Theorems
\ref{ord:canonical} and \ref{ord:threedirections}. Record the stronger
all-depth specialization in Theorem~\ref{ord:oneslope}.
\item Mark its tangential-curve question as answered by
Theorem~\ref{curve:main}, retaining the hypothesis that no prefix is
identically zero. The required jet order is a theorem with a sharpness
construction, not an empirical order suggested by examples.
\item Mark the cubic coincident-moment target as solved in specified
Tornheim coordinates by (C3)--(C6), and the digamma cube as reduced by
(C9). Keep the reduction of $\tau'''(0)$ to a chosen depth-one algebra
as an open question.
\item Cite the established Laurent work of Saha and of Matsumoto,
Onozuka and Wakabayashi; preserve the harmonic convention conversion
\eqref{harm:conversion}. Cite the classical triple-Hurwitz and
Mellin--Tornheim inputs when integrating the cubic section.
\item Distinguish the spatial finite part, the regularized
multivariate germ, and the constant on a specified spectral curve.
Formula (C10) gives their concrete cubic comparison. It should be
added as a normalization identity, not announced as an erratum to a
source that has not been shown to confuse them.
\item Preserve the conjectural status of the short $S_6$ and revised
$S_8$ candidates. Their existing proximity certificates and restricted
separation results are not proofs or general disproofs of the proposed
period identities.
\end{enumerate}

\subsection{Reproducible package}

The package contains the modular article source, a flattened
standalone TeX source, the compiled PDF, exact-check programs,
independent numerical programs, recorded JSON results, pinned source
hashes, and a claim-by-claim integration ledger. The proofs are
ordinary mathematical proofs; no Lean or other proof-assistant
formalization is claimed. The build script compiles the article
without fetching any mathematical data. Verification uses Python,
SymPy and mpmath; its commands and recorded versions are in the
package README.

The proposed placement follows the repository's existing thematic
research-report structure. The delivered section sources can also be
integrated into the canonical manuscript after notation reconciliation.
The present delivery does not modify the repository or promote its
unrelated conjectures.

% END sections/06_verification.tex


% BEGIN sections/07_questions.tex
\section{Further research questions and exact-identity programme}
\label{sec:questions}

The following questions begin where the proved formulas stop. They
are proposed research directions, not assertions disguised as
conjectures. Their intended outputs are exact identities, finite
coordinate descriptions, or proofs of specific candidate equalities.

\subsection{Independent slopes beyond depth three}

\paragraph{Q1. The first genuinely new ordered coefficients.}
At depth four the finite part of a general ray uses the linear jet
of $R_3$ and the quadratic jet of $R_2$, in addition to ordinary
Stieltjes data. Determine a complete finite basis of these
coefficients modulo the exact stuffle relations, shift identities,
and the known diagonal Gamma constraints. Then write a formula
analogous to \eqref{ord:threeFP} for all transverse quadruples of
slopes. The first objective is a proved spanning description; any
claim of arithmetic independence would require a separate argument.

\paragraph{Q2. Classify the directions with ordinary-zeta closure.}
The family with equal trailing slopes has ordinary Hurwitz closure
at every depth. At depth three the coefficient of the ordering
coordinate is $(c_2-c_3)/c_1$. At depth four and beyond, determine
the algebraic loci on which all genuinely nonsymmetric coefficient
contributions cancel. One can first work in the formal quotient by
proved identities, where the question becomes an exact finite
linear-algebra problem. Accidental numerical relations at a specific
shift must be distinguished from identities valid as functions of $a$.

\paragraph{Q3. Several blocks of equal slopes.}
Replace $(p,q,\ldots,q)$ by a sequence of several constant blocks.
Does each fixed number of blocks admit a generating function involving
a finite collection of lower-depth regular germs, uniformly in the
block lengths? The all-equal and one-distinguished-slope cases suggest
an induction by ordered blocks, but the crossing terms need explicit
control. A useful deliverable would be an exact two-block formula at
all depths and a proof of its coordinate requirements.

\subsection{Harmonic, polylogarithmic, and Gamma relations}

\paragraph{Q4. Rational-shift distribution of the ordering coordinate.}
The shift relation \eqref{harm:shift} is exact. The existing
root-of-unity filter in the incoming endpoint report already gives
finite coloured descriptions of rational-shift ordered sums and
their fixed spectral derivatives \cite{repoendpoint}. The harder
question is which character-weighted combinations of $\eta(r/q)$
eliminate the residual coloured ordering coefficients and reduce
entirely to ordinary Hurwitz or Dirichlet $L$-derivatives. Determine
the exact elimination spaces first at conductors $2,3,4,6$, then
seek a conductor-uniform rule. The inclusive harmonic literature
\cite{Kargin} is a starting point; the strict/inclusive conversion
must remain explicit.

\paragraph{Q5. Normalized parameter primitives.}
Equation~\eqref{harm:derivative} reduces $\eta'(a)$ to absolutely
convergent ordered sums. Seek a useful normalized antiderivative in
the shift parameter, expressed through multiple Gamma or Barnes-type
functions and spectral zeta derivatives. A candidate must specify
its additive constant and prove differentiation back to $\eta$.
The whole higher harmonic hierarchy is already represented by
\eqref{harm:allcoefficients}--\eqref{harm:allconversion}; seek a
parameter-primitive relation compatible with that hierarchy and with
the exact shift recurrence. Any regularized integral must carry its
finite normalization explicitly.

\paragraph{Q6. Logarithmic decorations of convergent primitive families.}
One Euler integration raises the outer order from one to two, so
$Z_{r+1}(2+u_0,1+u_1,\ldots,1+u_r;a)$ is holomorphic near the origin.
For this ordinary primitive the endpoint limit and every fixed
spectral derivative commute by domination; there is no unexplained
regularization defect. The incoming endpoint work already evaluates
the unweighted integer-shift primitive hierarchy
\cite{repoendpoint}. The next target is to classify which independent
logarithmic decorations, obtained by differentiating the separate
indices, still reduce to ordinary Gamma, zeta and Stieltjes jets.
Begin with the convergent depth-two derivatives appearing in
\eqref{harm:derivative}, derive their exact integral relations, and
compare them with the height-one Gamma generating relations.

\subsection{Tornheim derivatives and higher coincident products}

\paragraph{Q7. Evaluate the remaining cubic Tornheim derivative.}
Can $\tau'''(0)$ be reduced to a specified algebra generated by
$\gamma_m$, ordinary $\zeta^{(j)}(n)$, logarithms, and perhaps Barnes
Gamma constants? Formula (C9) turns this into an equivalent exact
evaluation question for the coordinate-normalized digamma cube.
The independent Mellin representation permits systematic analytic
manipulations and high-precision exploration, but an integer relation
would remain a candidate until a functional identity proves it.
Existing Tornheim derivative work and the recent Kronecker-limit
literature should be checked before any priority claim
\cite{BorweinDilcher,BaileyBorwein,Onodera2021,SS2026}.

\paragraph{Q8. Minimal coordinates for all cubic moments.}
Formula (C4) is a complete generating prescription. It does not yet
classify how many independent completed Tornheim jets are needed at
each total Stieltjes index $m_1+m_2+m_3$. Determine exact reductions
under permutation symmetry, integration by parts with the stated
finite-part convention, and Tornheim functional equations. The
first useful target is a table for total indices through three,
with proofs and a transparent distinction between a spanning family
and an arithmetically minimal one.

\paragraph{Q9. Four and more coincident factors.}
Extend the cubic Fourier kernel to the zero-frequency sums arising
from four Hurwitz factors. The endpoint resonances can be organized
by subsets of singular factors and their smooth Taylor degrees;
the genuinely new work is the global multiple-sum coordinate and its
functional reductions. Begin with a fixed-coordinate formula for
$\FP\int_0^1\psi(x)^4\,dx$, keeping all spectral-counterterm finite
contributions. The natural higher-depth function should be defined
independently, as (C6) defines $\tau$ here.

\subsection{Singular geometry and the Gaussian programme}

\paragraph{Q10. Restrictions lying inside a polar hyperplane.}
Theorem~\ref{curve:main} handles every analytic curve whose prefixes
are not identically zero. A curve contained in $P_j=0$ needs an
additional prescription. Compare residue-first, regular-germ-first,
and finite-part-in-a-normal-variable constructions, and derive their
exact differences. Decide which prescriptions respect the stuffle
product and how those choices propagate under changes of coordinates.
This should be formulated as a comparison theorem; no canonical
value follows just from substituting into a rational slope formula.

\paragraph{Q11. A functional proof of the short Gaussian candidates.}
The current $S_6$ and revised $S_8$ candidates require identities
between specified polylogarithmic periods. A productive next step
would be an exact functional relation, differential equation with
enough boundary data, or a finite symbolic certificate in a rigorously
defined iterated-integral algebra. Existing restricted separator
results constrain particular ansatz classes; they do not establish
general period independence. The ordered Hurwitz identities in this
article may supply useful coefficient relations, but no demonstrated
bridge presently converts them into a proof of either Gaussian
candidate. That bridge itself should be stated as a concrete research
target rather than assumed.

% END sections/07_questions.tex


% BEGIN references.tex
\begin{thebibliography}{99}
\addcontentsline{toc}{section}{References}

\bibitem{repomain}
V. Reshetnikov and the ProveIt research programme,
\emph{Polylogarithms and their Arithmetic Bridges}, canonical
manuscript and status records, snapshot \texttt{4c1286161c42},
10 October 2026; full commit identifier in Section~\ref{sec:scope}.
\href{https://github.com/VladimirReshetnikov/ProveIt/tree/4c1286161c425eca957b4b9f7cd436c37a4c98f3/Analysis/Polylogarithms/docs/manuscript}{Pinned manuscript directory}.

\bibitem{repojets}
ProveIt research continuation,
\emph{Resonant Gamma Jets and Directional Zeta Identities:
Discrete Stieltjes products, shifted polylogarithms, and local formulas
for higher spectral derivatives}, 10 October 2026.
Incoming archive \src{ProveIt_Resonant_Jets_2026-10-10.zip},
\src{sections/08-audit-research.tex}, further-research questions Q1 and Q3.
\href{https://github.com/VladimirReshetnikov/ProveIt/tree/4c1286161c425eca957b4b9f7cd436c37a4c98f3/docs/incoming}{Pinned incoming directory}.

\bibitem{repocoincident}
ProveIt research continuation,
\emph{Coincident-Point Stieltjes Calculus: Exact Quadratic Closure,
Polygamma Parity, and Polylogarithmic Collision Subtraction},
10 October 2026. Incoming archive
\src{ProveIt_Coincident_Stieltjes_2026-10-10.zip},
\src{article/article.tex}, subsection ``Cubic coincident moments and
higher-depth coordinates.''

\bibitem{repoendpoint}
ProveIt research continuation,
\emph{Endpoint Regularization of Logarithmic Harmonic Polylogarithms:
Stieltjes coefficients, Gamma transfer, spectral finite parts, and
exact antiderivatives}, 10 October 2026. Incoming archive
\src{ProveIt_Endpoint_Regularization_2026-10-10.zip},
\src{article/endpoint_regularization.tex}.

\bibitem{Saha}
B. Saha, Multiple Stieltjes constants and Laurent type expansion of the
multiple zeta functions at integer points,
\emph{Selecta Math. (N.S.)} \textbf{28} (2022), article 6.
\url{https://doi.org/10.1007/s00029-021-00719-1}.
Author preprint: \url{https://arxiv.org/abs/1902.04389}.

\bibitem{MOW}
K. Matsumoto, T. Onozuka and I. Wakabayashi,
Laurent series expansions of multiple zeta-functions of Euler--Zagier
type at integer points,
\emph{Math. Z.} \textbf{295} (2020), 623--642.
\url{https://doi.org/10.1007/s00209-019-02337-2}.
Author preprint: \url{https://arxiv.org/abs/1601.05918}.

\bibitem{Kargin}
L. Karg\i n, A. Dil, M. Cenkci and M. Can,
On the Stieltjes constants with respect to harmonic zeta functions,
\emph{J. Math. Anal. Appl.} \textbf{525} (2023), article 127302.
\url{https://doi.org/10.1016/j.jmaa.2023.127302}.
Author preprint: \url{https://arxiv.org/abs/2304.03517}.

\bibitem{EMTornheim}
O. Espinosa and V. H. Moll,
The evaluation of Tornheim double sums. Part 1,
\emph{J. Number Theory} \textbf{116} (2006), 200--229.
\url{https://doi.org/10.1016/j.jnt.2005.04.008}.
Author preprint: \url{https://arxiv.org/abs/math/0505647}.

\bibitem{BorweinDilcher}
J. M. Borwein and K. Dilcher,
Derivatives and fast evaluation of the Tornheim zeta function,
\emph{Ramanujan J.} \textbf{45} (2018), 413--432.
\url{https://doi.org/10.1007/s11139-017-9890-9}.

\bibitem{BaileyBorwein}
D. H. Bailey and J. M. Borwein,
Computation and experimental evaluation of Mordell--Tornheim--Witten
sum derivatives,
\emph{Exp. Math.} \textbf{27} (2018), 370--376.
\url{https://doi.org/10.1080/10586458.2017.1295687}.
Author preprint:
\url{https://www.davidhbailey.com/dhbpapers/omega-numerics.pdf}.

\bibitem{Onodera2021}
K. Onodera, On multiple Hurwitz zeta function of Mordell--Tornheim type,
\emph{Int. J. Number Theory} \textbf{17} (2021), 2327--2360.
\url{https://doi.org/10.1142/S1793042121500913}.

\bibitem{SS2026}
S. Sathyanarayana and N. G. Sharan,
Mordell--Tornheim zeta function: Kronecker limit type formulas and
special values,
\emph{Adv. Appl. Math.} \textbf{177} (2026), article 103075.
\url{https://doi.org/10.1016/j.aam.2026.103075}.
Author preprint: \url{https://arxiv.org/abs/2510.10093}.

\bibitem{DLMFHurwitz}
NIST Digital Library of Mathematical Functions,
Section 25.11, Hurwitz zeta function.
\url{https://dlmf.nist.gov/25.11}. Accessed 10 October 2026.

\bibitem{DLMFPolylog}
NIST Digital Library of Mathematical Functions,
Section 25.12, Polylogarithms.
\url{https://dlmf.nist.gov/25.12}. Accessed 10 October 2026.

\end{thebibliography}

% END references.tex

\end{document}
